\scanpage{63}
\deleted{$\alpha(D_1)\in\Gamma(S_0,L_0)$ définie par la condition $\beta(\alpha(D_1))=\gamma(b_0)-D_1$ (où $\beta:L_0\xrightarrow{\sim}\operatorname{Ker}(L_1\longrightarrow L_0)$ est l'isom. canonique induit par la multiplication par 2 dans $L_1$)}
\margin{\deleted{$D_{\unsure{0|1}}\equiv b_{\unsure{1|0}}^2\bmod 2$}}

que $D_1=D\unsure{\upharpoonright|\vert}S_1$ \inserted{$\in\Gamma(S_1,L_1^{\otimes2})$} est \uline{localmt} le carré d'une section de $L_1$. \inserted{Il suffit, et il faut} si $H^1(S_0,L_0)=0$, \deleted{cette condition équivaut au} que \deleted{\unsure{$D_1$}} il existe $b_1\in\Gamma(S_1,L_1)$ tel que
\[b_1^2=D_1.\]

Enfin, il suffit, et il faut si $H^1(S,L)\unsure{-|\to}H^1(S\setminus S_0,L)$ est injectif, qu'il existe $b\in\Gamma(S,\underline L)$ et $c\in\Gamma(S,\underline L^{\otimes\unsure{2}})$ tels qu'on ait
\[b^2-4c=D.\]

En effet, la première assertion résulte du fait que si $S_0$ est réduit, la section $D_0=D_{S_0}\deleted{\unsure{|S_0}}$ admet au plus une \uline{section}. Le reste suit de façon évidente.

\uline{NB} La condition que $D_1=b_1^2$ localement se décompose en deux

a) $D_0=b_0^2$ (où $b_0\in\Gamma(S_0,L_0)$ est uniquement déterminé \inserted{i.e. \deleted{\unsure{$D$}} la classe $\unsure{\partial|\delta}(b_0)\in H^1(S_0,\mu_2)$ est nulle !}

b) $D_1=\gamma(b_0)$, ce qu'on peut aussi écrire sous la forme
\[\alpha(D_1)=0\]
où $\alpha(D_1)\in\Gamma(S_0,L_0)$ est défini par la condition
\[\beta(\alpha(D_1))=\gamma(b_0)-D_1\]
où $\beta:\deleted{\unsure{\operatorname{Ker} S_0}}L_0\unsure{\}}\xrightarrow{\sim}\operatorname{Ker}(L_1\longrightarrow L_0)$
% archive sheet number: 62
